\scanpage{206}
$\deleted{\unsure{\Gamma}}G$ \quad $k$-gpe loc. de type fini sur $k$,
$G_0$ \quad composante connexe
\[
\Gamma=G/G_0
\]
\[
\Gamma\xrightarrow{\pi_n}\Gamma\qquad {}_n\Gamma=\pi_n^{-1}(e)
\]
c'est un sous-schéma \uline{ouvert} et \uline{fermé} [car $\Gamma$
discret \uline{sép}] ens des $a$ tels que $\pi_n(a)=e$

(i) La suite des $\,{}_n\Gamma$ est loc. stationnaire, \deleted{\unsure{or}} (\unsure{trivial}
car $\Gamma$ est discret] \deleted{\unsure{alors}} son
sup est un sous-schéma \uline{ouvert} \uline{induit} $\Gamma^\tau$
de $\Gamma$ [N.B. c'est vrai pour tt $k$-groupe
loc. fini sur $k$, pas néc. séparable, mais
c'est moins évident, le cas radiciel
demande une démonstration spéciale]

\deleted{\unsure{(ii)}} Soit $a\in\Gamma$, alors $a\in{}_n\Gamma\deleted{\unsure{^\tau}}$ sss \deleted{\unsure{$a=e$}}
\deleted{il existe $n$}\inserted{$\pi_n(a)=e$}
$a\in\Gamma^\tau\Longleftrightarrow\exists n,\ \pi_n(a)=e\Longleftrightarrow$ il existe un sous-
\deleted{tel que $a$}
\deleted{(ii) Soit} groupe \deleted{ouvert} [néc. ouvert car $\Gamma$ sép sur $k$]
$\Gamma'$ \uline{fini} sur $k$ tel que $a\in\Gamma'\Longleftrightarrow$
l'un des $a^n$ est \uline{fermé}.

(iii) Si $\Gamma$ est commutatif, alors $\Gamma^\tau$ est le
sup des sous-groupes [ouverts] \uline{finis}
de $\Gamma$. C'est aussi le plus petit sous-groupe
tel que $\Gamma/\Gamma^\tau$ soit ``sans torsion''

(iv) Soit $G^\tau$ l'image inverse de $\Gamma^\tau$. \deleted{\unsure{Alors}}\inserted{C'est}
un sous-schéma ouvert de $G$, et
$a\in G^\tau\Longleftrightarrow\exists n,\ \pi_n(a)\in G_0\Longleftrightarrow$ Il existe
un sous-groupe ouvert $H$ de $G$ contenant \unsure{$a$|les $a$} \deleted{\unsure{$a^n$|les}}
$\Longleftrightarrow\exists$ \unsure{ss-ensemble|ss-schéma} \uline{de type fini} de $G$ contenant les $a^n$
% archive sheet number: 205
